% Use the existing owner-derived source mechanism; do not edit frozen originals.
% Source-only commented blocks are deliberately excluded from TeX macro arguments.
% Apply all exact byte-preserving edits from source-transform.json before tokenization.
\OLSADeclareDocumentTextCorrection{OLP-0037}{FA0037-TWO-DIRECTIONS}
{مقایسه را به دو حالت تجزیه کنیم}
{مقایسه را در دو جهت جداگانه انجام دهیم}
\OLSADeclareDocumentTextCorrection{OLP-0038}{FA0038-EMPTY-INITIAL-SEGMENT}
{!!{element}sی~$A$ است.}
{!!{element}sی~$A$ است.
\emph{یادداشت دربارهٔ حالت تهی:} در این جمله، قطعهٔ آغازینِ متناظر با
مجموعهٔ تهی نیز تهی در نظر گرفته می‌شود. تعریفِ بعدی، شمارابودنِ مجموعهٔ
تهی را جداگانه تصریح می‌کند.}
\OLSADeclareDocumentTextCorrection{OLP-0038}{FA0038-CEILING-PRECISION}
{که $x$ را رو به بالا
تا نزدیک‌ترین عدد صحیح گرد می‌کند.}
{که $x$ را به کوچک‌ترین عدد صحیحِ بزرگ‌تر یا مساویِ آن می‌برد.}
\OLSADeclareDocumentTextCorrection{OLP-0039}{FA0039-INFINITE-QUALIFIER}
{این مطلب شاید شگفت‌آور باشد. زیرا گفتن اینکه $A$}
{این مطلب شاید شگفت‌آور باشد. برای یک مجموعهٔ نامتناهی، گفتن اینکه $A$}
\OLSADeclareDocumentTextCorrection{OLP-0039}{FA0039-CONVENTION-DISCLOSURE}
{«بیشتر» است.}
{«بیشتر» است.
\emph{یادداشت ویراستاری:} نبودنِ تناظر دوسویی با کل اعداد طبیعی،
مجموعه‌های متناهی را نیز شامل می‌شود؛ قید نامتناهی در این توضیح ضروری است.
در مورد مجموعه‌های دلخواه، تعبیرِ مقایسهٔ اندازه‌ها در اینجا با فرض اصل
انتخاب خوانده می‌شود. برهان‌های قطریِ زیر به این فرض نیاز ندارند.}
\OLSADeclareDocumentTextCorrection{OLP-0039}{FA0039-METHOD-INFINITE-SCOPE}
{برای اثبات اینکه مجموعه‌ای !!{nonenumerable} است،}
{برای اثبات اینکه مجموعه‌ای نامتناهی !!{nonenumerable} است،}
\OLSADeclareDocumentTextCorrection{OLP-0039}{FA0039-INFINITE-LIST-PREMISE}
{هر برشماری از یک زیرمجموعهٔ $\Bin^\omega$ را در نظر بگیرید.}
{هر فهرستِ نامتناهی از اعضای $\Bin^\omega$ را در نظر بگیرید.}
\OLSADeclareDocumentTextCorrection{OLP-0039}{FA0039-COORDINATE-ORDER}
{$s_n(m)$ رقم $n$اُمِ رشتهٔ $m$اُمِ}
{$s_n(m)$ رقم $m$اُمِ رشتهٔ $n$اُمِ}
\OLSADeclareDocumentTextCorrection{OLP-0039}{FA0039-COORDINATE-DISCLOSURE}
{در آن $s_n(m)$ در سطر $n$اُم و ستون $m$اُم قرار گرفته است:}
{در آن $s_n(m)$ در سطر $n$اُم و ستون $m$اُم قرار گرفته است.
\emph{یادداشت ویراستاری:} ترتیب شاخص‌های رقم و رشته در جملهٔ متن مبدأ
جابه‌جا شده بود؛ اینجا با آرایه هماهنگ شده است:}
\OLSADeclareDocumentTextCorrection{OLP-0039}{FA0039-FLIP-BOTH-WAYS}
{$1$ را به $0$ و هر $1$ را به~$0$}
{$1$ را به $0$ و هر $0$ را به~$1$}
\OLSADeclareDocumentTextCorrection{OLP-0039}{FA0039-FLIP-DISCLOSURE}
{تغییر می‌دهیم. به‌بیان انتزاعی‌تر،}
{تغییر می‌دهیم.
\emph{یادداشت ویراستاری:} جهتِ دومِ تبدیل در جملهٔ متن مبدأ به‌اشتباه
تکرار شده بود؛ عبارت با تعریفِ زیر هماهنگ شده است. به‌بیان انتزاعی‌تر،}
\OLSADeclareDocumentTextCorrection{OLP-0039}{FA0039-LIST-CONCLUSION}
{نشان داده‌ایم که هر برشماریِ دلخواه از زیرمجموعه‌ای از $\Bin^\omega$،}
{نشان داده‌ایم که هر فهرستِ نامتناهیِ دلخواه از اعضای $\Bin^\omega$،}
\OLSADeclareDocumentTextCorrection{OLP-0039}{FA0039-NONFINITE-DISCLOSURE}
{از مجموعهٔ $\Bin^\omega$ وجود ندارد؛ یعنی $\Bin^\omega$
!!{nonenumerable} است.}
{از مجموعهٔ $\Bin^\omega$ وجود ندارد؛ یعنی $\Bin^\omega$
!!{nonenumerable} است.
\emph{توضیح:} خودِ مجموعهٔ رشته‌های دودوییِ نامتناهی، نامتناهی است؛ پس
برشماریِ فرضیِ کل آن نمی‌تواند دامنه‌ای متناهی داشته باشد. به همین دلیل
بررسیِ فهرست‌های نامتناهی در این برهان کافی است.}
\OLSADeclareDocumentTextCorrection{OLP-0039}{FA0039-CONDITIONAL-FOOTNOTE}
{\oliflabeldef{sfr:card-arithmetic:card-opps:sec}{\footnote{دقیق‌تر
آن است که مقرر کنیم $\Bin^\omega$ مجموعهٔ همهٔ دنباله‌های $\omega$تایی
از $0$ها و $1$ها، یعنی مجموعهٔ
$\funfromto{\omega}{\{0,1\}}$، است. اما معنای این مطلب تنها در
\olref[sfr][card-arithmetic][card-opps]{sec} روشن خواهد شد.}{} بااین‌حال،
این صورت‌بندیِ اندکی مسامحه‌آمیز فعلاً نباید موجب هیچ ابهامی شود.}}
{\oliflabeldef{sfr:card-arithmetic:card-opps:sec}{\footnote{دقیق‌تر
آن است که مقرر کنیم $\Bin^\omega$ مجموعهٔ همهٔ دنباله‌های $\omega$تایی
از $0$ها و $1$ها، یعنی مجموعهٔ
$\funfromto{\omega}{\{0,1\}}$، است. اما معنای این مطلب تنها در
\olref[sfr][card-arithmetic][card-opps]{sec} روشن خواهد شد.}}{}
بااین‌حال، این صورت‌بندیِ اندکی مسامحه‌آمیز فعلاً نباید موجب هیچ ابهامی شود.}
\OLSADeclareDocumentTextCorrection{OLP-0039}{FA0039-MATRIX-MEMBERSHIPS-ROW0}
{N_0 = \{ & \mathbf{0}, & 1, & 2, & & & & \dots\}\\}
{N_0 = \{ & \mathbf{0}, & 1, & 2, & 3, & 4, & 5, & \dots\}\\}
\OLSADeclareDocumentTextCorrection{OLP-0039}{FA0039-MATRIX-MEMBERSHIPS-ROW3}
{N_3 = \{ &  &  & 2, & \mathbf{3}, & 4, & & \dots\}\\}
{N_3 = \{ &  &  & 2, & \mathbf{3}, & 4, & 5, & \dots\}\\}
\OLSADeclareDocumentTextCorrection{OLP-0039}{FA0039-MATRIX-DISCLOSURE}
{باشند؛ آنگاه آرایهٔ ما چنین آغاز می‌شود:}
{باشند؛ آنگاه آرایهٔ ما چنین آغاز می‌شود.
\emph{یادداشت ویراستاری:} در آرایهٔ متن مبدأ، سه عضوِ سطر نخست و یک
عضوِ سطر چهارم از قلم افتاده بود. در اینجا همهٔ خانه‌ها مطابق عضویت در
مجموعه‌های داده‌شده تکمیل شده‌اند:}
\OLSADeclareDocumentTextCorrection{OLP-0040}{FA0040-SURJECTIVE-LIST-CONVENTION}
{برشمارد، آنگاه $f(x_1)$، $f(x_2)$، \dots{} مجموعهٔ~$B$ را برمی‌شمارد.}
{برشمارد، آنگاه $f(x_1)$، $f(x_2)$، \dots{} هر عضوِ مجموعهٔ~$B$ را
دست‌کم یک بار در بر خواهد داشت.
\emph{یادداشت دربارهٔ تعریف:} این فهرست ممکن است تکرار داشته باشد و
بنابراین هنوز لزوماً برشماری به معنای تناظر دوسویی نیست. با نگه‌داشتن
نخستین رخدادِ هر مقدار و شماره‌گذاریِ دوباره، برشماریِ دوسوییِ لازم به
دست می‌آید؛ حالتِ تهی جداگانه طبق تعریف پوشش داده می‌شود.}
\OLSADeclareDocumentTextCorrection{OLP-0040}{FA0040-CHARACTERISTIC-SYMBOL}
{$s_{k}$ باشد}
{$s$ باشد}
\OLSADeclareDocumentTextCorrection{OLP-0040}{FA0040-CHARACTERISTIC-ONE}
{$s_{k}(n) = 1$}
{$s(n) = 1$}
\OLSADeclareDocumentTextCorrection{OLP-0040}{FA0040-CHARACTERISTIC-ZERO}
{$s_k(n) = 0$}
{$s(n) = 0$}
\OLSADeclareDocumentTextCorrection{OLP-0040}{FA0040-SYMBOL-DISCLOSURE}
{اگر $n \in N$، و در غیر این صورت $s(n) = 0$.}
{اگر $n \in N$، و در غیر این صورت $s(n) = 0$.
\emph{یادداشت ویراستاری:} زیرنویسِ تعریف‌نشدهٔ متن مبدأ حذف شده است؛
نماد رشته در این تعریف به همان مجموعهٔ ورودی وابسته است.}
\OLSADeclareDocumentTextCorrection{OLP-0040}{FA0040-BIJECTIVE-SPECIFIC-LIST}
{فهرست باید همهٔ~$\Bin^\omega$ را برشمارد.}
{فهرست باید همهٔ~$\Bin^\omega$ را برشمارد.
این نگاشت یک‌به‌یک نیز هست، زیرا رشتهٔ حاصل، عضویتِ هر عدد طبیعی در
مجموعهٔ ورودی را تعیین می‌کند. پس در این مثال، فهرستِ حاصل از یک
برشماریِ دوسویی، خود نیز بدون تکرار است.}
\OLSADeclareSetup{OLP-0037}{\olsa_source_correction_install_document_text:}
\OLSADeclareSetup{OLP-0038}{\olsa_source_correction_install_document_text:}
\OLSADeclareSetup{OLP-0039}{\olsa_source_correction_install_document_text:}
\OLSADeclareSetup{OLP-0040}{\olsa_source_correction_install_document_text:}
